% =====================================================================
%   附录 G：家务整理 (House Holding)
% =====================================================================

\section{家务整理}\label{app:hh}

\subsection{数据集细节}

\subsubsection*{构造细节}

ALFWorld 基准由一系列模拟家务场景的文本环境组成，提供了一个智能体可借助
文本界面执行决策任务的交互式环境。给定家务环境描述与目标指令，智能体的目标
是将复杂的高层目标分解为一系列简洁的动作。每一步之后，智能体会收到来自环境
的反馈，使其能够动态调整计划，并转向后续任务，最终达成主目标。

ALFWorld 数据集中每条评估样本包含以下内容：
\begin{itemize}
    \item \textbf{环境描述（Environment Description）}：包括智能体初始位置
          以及所在房间中物体及其 ID 的快照，对整体家务环境的详细描述。
    \item \textbf{目标（Objective）}：智能体在该环境中需达成的目标，通常
          需要多步推理与探索（例如：将台灯放到桌上）。
    \item \textbf{仿真环境（Simulated Environment）}：在智能体的每一步动作
          之后，仿真环境会提供即时反馈，并评估智能体是否已完成该任务。
\end{itemize}

在数据集中，我们使用了 ALFWorld 数据集中评估分布外（eval out of distribution）
划分中 134 个可解的问题。所有问题被分为 6 类：取放（pick and place）、
取放清洁后置（pick clean then place）、取放加热后置（pick heat then place）、
取放冷却后置（pick cool then place）、观察物品（look at obj）与取两件物品
（pick two obj）。

\subsubsection*{评测设置}

由于问题本身的固有复杂性以及对输出格式的高标准要求，我们采用 1-shot 评测
设置：对每个类别的问题，我们从训练集中选取同一类别的一个相对简单且完整的
交互过程作为示例。沿用 ReAct \citep{yao2023breact} 的做法，我们采用其开源仓库\footnote{\url{https://github.com/ysymyth/ReAct}}
中的少样本示例与提示词。此外，若 LLM 输出的格式无效，我们使用 BLEU 指标来
评估输出与所有合法动作选项之间的相似度，并将相似度最高的选项作为模型
本轮的动作。

对每条样本，评估过程可分为 2 个部分：
\begin{enumerate}
    \item \textbf{初始化}：我们向模型描述任务并提供一个成功的示例。随后
          我们详细描述环境，并阐明需要完成的目标。
    \item \textbf{交互}：模型根据前序交互反馈与环境信息，生成若干思考
          与下一动作。在收到模型的动作后，环境给出反馈（环境状态的变化
          或模型所观察到的信息）。该过程反复进行，直到模型成功达成目标
          （视为成功）或达到最大动作次数（视为失败）。值得注意的是，在
          多次未果之后，模型有时会反复输出相同内容。为节省评测时间，
          我们规定：若模型连续三次输出相同内容，则视为因重复而失败。
\end{enumerate}

\subsubsection*{评测指标}

我们以总体成功率（Success Rate）作为模型性能的度量，即模型成功完成任务的
数量与任务总数之比。

\subsection{提示词示例}

为使输出格式与仿真环境所支持的合法命令对齐，我们采用 1-shot 评测设置，
即在指令后拼接一个成功完成任务的示例。在交互开始时，我们使用如下指令
向模型描述任务：
\begin{quote}\small\ttfamily
Interact with a household to solve a task. Imagine you are an intelligent
agent in a household environment and your target is to perform actions to
complete the task goal. At the beginning of your interactions, you will
be given the detailed description of the current environment and your
goal to accomplish. For each of your turn, you will be given a list of
actions which you can choose one to perform in this turn. You should
choose from two actions: ``THOUGHT'' or ``ACTION''. If you choose
``THOUGHT'', you should first think about the current condition and plan
for your future actions, and then output your action in this turn. Your
output must strictly follow this format:
``THOUGHT: your thoughts. \textbackslash n ACTION: your next action\textbackslash n'';
If you choose ``ACTION'', you should directly output the action in this
turn. Your output must strictly follow this format:
``ACTION: your next action\textbackslash n''. After your each turn, the
environment will give you immediate feedback based on which you plan
your next few steps. If the environment output ``Nothing happened'',
that means the previous action is invalid and you should try more
options.\\
\\
Reminder:
\begin{enumerate}[label=\arabic*.,leftmargin=*]
    \item the action must be chosen from the given available actions.
          Any actions except provided available actions will be regarded
          as illegal.
    \item Think when necessary, try to act directly more in the process.
\end{enumerate}
\end{quote}

数据集中的所有任务共分为 6 类。为更好地引导模型完成目标，我们从每个类别
中选取一个相对简单且成功完成的样例作为 1-shot 示例。下面给出一个示例
（对应"取放清洁后置"类任务；提示词详见原论文）：

\begin{quote}\small\ttfamily
User: You are in the middle of a room. Looking quickly around you, you
see a cabinet 4, a cabinet 3, a cabinet 2, a cabinet 1, a countertop 1,
a garbagecan 1, a handtowelholder 2, a handtowelholder 1, a sinkbasin 2,
a sinkbasin 1, a toilet 1, a toiletpaperhanger 1, and a towelholder 1.
Your task is to: put some spraybottle on toilet.\\
\\
Agent: THOUGHT: The task is to put some spraybottle on toilet, so first
I need to search the room for some spraybottle. After finding the
spraybottle, I need to take it to the toilet. I will start my search on
cabinets. \textbackslash n ACTION: go to cabinet 1\\
User: On the cabinet 1, you see a cloth 1, a soapbar 1, a soapbottle 1.\\
Agent: ACTION: go to cabinet 2\\
User: The cabinet 2 is closed.\\
Agent: ACTION: open cabinet 2\\
User: You open the cabinet 2. The cabinet 2 is open. In it, you see a
candle 1, and a spraybottle 2.\\
Agent: ACTION: take spraybottle 2 from cabinet 2\\
User: You pick up the spraybottle 2 from the cabinet 2.\\
Agent: ACTION: go to toilet 1\\
User: On the toilet 1, you see a soapbottle 2.\\
Agent: ACTION: put spraybottle 2 in/on toilet 1\\
User: You put the spraybottle 2 in/on the toilet 1.
\end{quote}
